% =====================================================================
%  Pseudo-code for the simulation and the estimation algorithms
%  (Appendix B.3 of "Spatial Comovements").
%
%  This file is a FRAGMENT: % =====================================================================
%  Pseudo-code for the simulation and the estimation algorithms
%  (Appendix B.3 of "Spatial Comovements").
%
%  This file is a FRAGMENT: % =====================================================================
%  Pseudo-code for the simulation and the estimation algorithms
%  (Appendix B.3 of "Spatial Comovements").
%
%  This file is a FRAGMENT: % =====================================================================
%  Pseudo-code for the simulation and the estimation algorithms
%  (Appendix B.3 of "Spatial Comovements").
%
%  This file is a FRAGMENT: \input{pseudocode_body} it where you want the
%  algorithms to appear (e.g. right after Appendix B.3.2).
%
%  Required packages in the master preamble:
%      \usepackage{amsmath,amssymb}
%      \usepackage{algorithm}
%      \usepackage{algpseudocode}
%      \usepackage{float}          % for the [H] float placement used below
%
%  Cross-references: equation numbers are written literally as (B12), (B22),
%  (B31), (B32) to match the current draft. Replace them by \eqref{...} once
%  those equations carry labels.
%  Recommended (renames Require/Ensure to Input/Output):
%      \renewcommand{\algorithmicrequire}{\textbf{Input:}}
%      \renewcommand{\algorithmicensure}{\textbf{Output:}}
%
%  `pseudocode.tex` in the same folder is a standalone driver that compiles
%  this file on its own, for checking.
% =====================================================================

\subsection{Pseudo-code}
\label{app:pseudocode}

This section states the two algorithms of Appendix~B.3 in
explicit form: Algorithm~\ref{alg:simulate} evaluates the model at a candidate
parameter vector $\mu$ and returns the simulated moment vector
$\widetilde{M}(\mu)$, and Algorithm~\ref{alg:smm} minimises the SMM criterion
over $\mu$ by particle swarm optimisation. Algorithms~\ref{alg:invertT},
\ref{alg:countN} and \ref{alg:pso} are the three subroutines they call: the
comparative-advantage inversion, the variety-count calibration, and the swarm
itself.

Throughout, $\mu=\bigl(\Omega_L,\{\Omega_s\},\alpha,\{A_r\}\bigr)$ is the vector
searched over; $a(r')$ denotes the attraction area of region $r'$;
$\mathcal{R}_a$ the regions of area $a$; $\mathcal{R}_D$ the set of downstream
buyers, $R_D=|\mathcal{R}_D|$; and $\mathcal{R}^+_s$ the regions inside active
attraction areas for sector $s$. The calibrated objects
\[
  \bigl\{\theta,\ \varepsilon,\ \lambda,\ \nu,\ \{\nu_s\},\
  \{w_{r's}\},\ \{\delta_r\},\ \{d_{r'r}\}\bigr\}
\]
are fixed once and for all and are suppressed from the argument lists.

% ---------------------------------------------------------------------
%  Algorithm 1 -- simulation
% ---------------------------------------------------------------------
\begin{algorithm}[H]
\caption{Simulation of the economy at a candidate $\mu$ (one criterion evaluation)}
\label{alg:simulate}
\begin{algorithmic}[1]
\Require candidate $\mu=(\Omega_L,\{\Omega_s\},\alpha,\{A_r\})$; empirical
  sourcing shares $\widehat{\gamma}_{as}$; empirical count moments
  $\widehat{G}_s(0)$ and bounds $[N^{lo}_s,N^{hi}_s]$; a fixed array of uniform
  draws $u_{r'\rho s}\in(0,1)$, $\rho=1,\dots,m$
\Ensure simulated moment vector $\widetilde{M}(\mu)$, and the calibrated
  $\{T_{as}\},\{\widehat{N}_s\}$
\vspace{2pt}
\State $\tau_{r'r}\gets d_{r'r}^{\,\alpha}$
  \Comment{trade costs; the only place $\alpha$ enters}
\vspace{2pt}
\State \textbf{Step 1 (comparative advantage).}
  $\{T_{as}\}\gets\Call{InvertT}{\alpha,\Omega_L,\{\Omega_s\},\{A_r\};\widehat{\gamma}}$
  \Comment{Algorithm~\ref{alg:invertT}}
\vspace{2pt}
\State \textbf{Step 2 (Ricardian competition).}
\For{$s=1,\dots,S$ \textbf{and} $\rho=1,\dots,m$}
    \For{$r'=1,\dots,R$}
        \State $z_{r'\rho s}\gets T_{a(r')s}^{1/\theta}\bigl(-\log u_{r'\rho s}\bigr)^{-1/\theta}$
          \Comment{Fr\'echet inverse c.d.f.}
    \EndFor
    \For{$r\in\mathcal{R}_D$}
        \State $r^\star(\rho,s,r)\gets\arg\min_{r'}\; w_{r's}\tau_{r'r}/z_{r'\rho s}$
          \Comment{winner of variety $\rho$ in market $(r,s)$}
        \State $p_{\rho s r}\gets w_{r^\star s}\tau_{r^\star r}/z_{r^\star\rho s}$
    \EndFor
\EndFor
\State $k_{r's}\gets\#\{\rho:\ r'=r^\star(\rho,s,r)\ \text{for some }r\in\mathcal{R}_D\}$,
  \quad $p_{r's}\gets k_{r's}/m$
  \Comment{win-anywhere counts and probabilities}
\vspace{2pt}
\State \textbf{Prices, sales and flows.} For every $r\in\mathcal{R}_D$:
\State \quad $P_{rs}\gets\bigl[\tfrac1m\sum_{\rho}p_{\rho s r}^{\,1-\nu_s}\bigr]^{1/(1-\nu_s)}$,\quad
  $P_r\gets\bigl[\sum_s\Omega_s P_{rs}^{1-\nu}\bigr]^{1/(1-\nu)}$
\State \quad $c_r\gets\bigl[\Omega_L w_r^{1-\lambda}+(1-\Omega_L)P_r^{1-\lambda}\bigr]^{1/(1-\lambda)}$,\quad
  $\tilde c_r\gets c_r/A_r$
\State \quad $y_r\gets \tilde c_r^{\,\varepsilon}\,\delta_r\big/\sum_{r''}\tilde c_{r''}^{\,\varepsilon}\delta_{r''}$,\quad
  $\pi_r\gets y_r$
  \Comment{iso-elastic demand, $E$ normalised to $1$}
\State \quad $X_{r'rs}\gets$ expenditure of $r$ on the varieties won by $r'$,
  from~(B12)
\vspace{2pt}
\State \textbf{Step 3 (variety count).}
  $\{\widehat{N}_s\}\gets\Call{CalibrateN}{\{k_{r's}\},m;\widehat{G}_s(0),[N^{lo}_s,N^{hi}_s]}$
  \Comment{Algorithm~\ref{alg:countN}}
\vspace{2pt}
\State \textbf{Step 4 (moments).} Assemble $\widetilde{M}(\mu)$ from the single
  solve above:
\State \quad (1) $\widetilde{\Omega}_L$, (2) $\widetilde{\Omega}_s$, (3) $\pi_r$
  \Comment{value blocks, from $\{X_{r'rs}\},\{c_r\},\{y_r\}$}
\State \quad (4) $\beta$: cloglog of ``does not supply'' on $\log d_{r'a(r')}$ and
  $\log z$, area${}\times{}$sector fixed effects, run on the simulated firm
  population $\{(r',\rho,s)\}$
\State \quad (5) $\widetilde{\gamma}_{as}\gets\sum_{r'\in\mathcal{R}_a}\sum_r X_{r'rs}\big/\sum_{r''}\sum_r X_{r''rs}$
\State \quad (6) $\widetilde{G}_s(\widehat{N}_s)$
  \Comment{count moment, evaluated at the calibrated $\widehat{N}_s$}
\State \Return $\widetilde{M}(\mu)$
\end{algorithmic}
\end{algorithm}

Three features of Algorithm~\ref{alg:simulate} are worth emphasising. First, the
arrows of~(B22) run one way: $T$ is calibrated without knowing
$N$, and $N$ does not revise $T$, so Steps~1--3 form a recursion rather
than a simultaneous system. Second, the simulation size $m$ is a Monte Carlo
device and is unrelated to $\widehat{N}_s$: the winner of a variety depends only
on relative production costs, so $p_{r's}$ can be estimated at any $m$ large
enough. Third, the uniform draws $u$ are held fixed across every evaluation of
the criterion within an optimisation stage (common random numbers), which is
what makes the criterion a deterministic function of $\mu$ that a swarm can
descend; they are scrambled Sobol sequences, constructed one low-dimensional
sequence per sector with regions as dimensions and varieties as points.
Independent pseudo-random draws are used instead when computing $\Sigma_{sim}$
and the numerical derivatives, where independence across replications is
required.

% ---------------------------------------------------------------------
%  Algorithm 2 -- comparative advantage inversion
% ---------------------------------------------------------------------
\begin{algorithm}[H]
\caption{\textsc{InvertT}: comparative advantage by matrix scaling with
general-equilibrium expenditure weights}
\label{alg:invertT}
\begin{algorithmic}[1]
\Require $\alpha,\Omega_L,\{\Omega_s\},\{A_r\}$; targets $\widehat{\gamma}_{as}$;
  damping $\varsigma\in(0,1]$; tolerance $\epsilon_T$; starting values $T^{(0)}_{as}$
\Ensure $\{T_{as}\}$ reproducing $\widehat{\gamma}_{as}$ exactly
\State $T\gets T^{(0)}$; normalise $T_{as}\gets T_{as}/T_{a_s^{\text{ref}}s}$ for each $s$
  \Comment{one reference area per sector fixes the scale}
\For{$t=1,2,\dots$ \textbf{until} convergence}
    \State $B_{ar}\gets\sum_{r'\in\mathcal{R}_a}(w_{r's}\tau_{r'r})^{-\theta}$,\quad
      $\Phi_{rs}\gets\sum_{a'}T_{a's}B_{a'r}$
    \State $\omega_{rs}\gets$ expenditure weights implied by $T$
      \Comment{price and sales block of Algorithm~\ref{alg:simulate}}
    \State $\gamma_{as}(T)\gets\sum_r\omega_{rs}\,T_{as}B_{ar}/\Phi_{rs}$
    \State $T^{+}_{as}\gets T_{as}\bigl(\widehat{\gamma}_{as}/\gamma_{as}(T)\bigr)^{\varsigma}$
      \Comment{multiplicative Sinkhorn update, in logs}
    \State $T^{+}_{as}\gets T^{+}_{as}/T^{+}_{a^{\text{ref}}_s s}$ for each $s$
    \State $\text{resid}\gets\max_{a,s}\bigl|\log T^{+}_{as}-\log T_{as}\bigr|$;\quad $T\gets T^{+}$
    \State \textbf{if} $\text{resid}<\epsilon_T$ \textbf{then break}
\EndFor
\State \Return $T$
\end{algorithmic}
\end{algorithm}

With expenditure weights held fixed, the multiplicative update is a Sinkhorn scaling of the strictly positive matrix
$B$, hence a contraction with a fixed point unique up to the sector-specific
scale removed by the reference-area normalisation. Recomputing $\omega_{rs}$
inside the loop adds the general-equilibrium response of prices and sales, so
the full map is a fixed-point iteration; the damping $\varsigma$ ($0.9$ in the
implementation) and the stopping rule on $\max_{a,s}|\Delta\log T_{as}|$ are what
the text refers to as requiring both sourcing shares and comparative advantages
to converge.

% ---------------------------------------------------------------------
%  Algorithm 3 -- variety count
% ---------------------------------------------------------------------
\begin{algorithm}[H]
\caption{\textsc{CalibrateN}: variety count by monotone integer search}
\label{alg:countN}
\begin{algorithmic}[1]
\Require win counts $k_{r's}$ from $m$ simulated varieties; targets
  $\widehat{G}_s(0)$; bounds $[N^{lo}_s,N^{hi}_s]$
\Ensure $\{\widehat{N}_s\}$
\For{$s=1,\dots,S$}
    \State define
      $\widetilde{G}_s(n)\gets\dfrac{1}{|\mathcal{R}^+_s|}\displaystyle\sum_{r'\in\mathcal{R}^+_s}
      \binom{m-k_{r's}}{n}\Big/\binom{m}{n}$
      \Comment{unbiased; decreasing in $n$}
    \If{$\widetilde{G}_s(N^{lo}_s)\le\widehat{G}_s(0)$}
        \State $\widehat{N}_s\gets N^{lo}_s$ \Comment{clamped at the lower bound}
    \ElsIf{$\widetilde{G}_s(N^{hi}_s)\ge\widehat{G}_s(0)$}
        \State $\widehat{N}_s\gets N^{hi}_s$ \Comment{clamped at the upper bound}
    \Else
        \State $(\underline{n},\overline{n})\gets(N^{lo}_s,N^{hi}_s)$
        \While{$\overline{n}-\underline{n}>1$}
            \State $n\gets\lfloor(\underline{n}+\overline{n})/2\rfloor$
            \State \textbf{if} $\widetilde{G}_s(n)>\widehat{G}_s(0)$ \textbf{then}
              $\underline{n}\gets n$ \textbf{else} $\overline{n}\gets n$
        \EndWhile
        \State $\widehat{N}_s\gets\arg\min_{n\in\{\underline{n},\overline{n}\}}
          \bigl|\widetilde{G}_s(n)-\widehat{G}_s(0)\bigr|$
    \EndIf
\EndFor
\State \Return $\{\widehat{N}_s\}$
\end{algorithmic}
\end{algorithm}

Because $\widetilde{G}_s$ is decreasing in $n$ and is evaluated in closed form
from the winning matrix of Step~2, the search costs $O(\log(N^{hi}_s-N^{lo}_s))$
evaluations and never re-simulates the Ricardian competition. The count is an
integer at every evaluation: it is never relaxed to the real line and never
rounded. A bound that binds is recorded and reported rather than silently
absorbed.

% ---------------------------------------------------------------------
%  Algorithm 4 -- PSO
% ---------------------------------------------------------------------
\begin{algorithm}[H]
\caption{\textsc{Pso}: particle swarm minimisation of $\mathcal{Q}$ over a box}
\label{alg:pso}
\begin{algorithmic}[1]
\Require objective $\mathcal{Q}$; box $[\ell,u]\subset\mathbb{R}^{d}$; swarm size
  $P$; iterations $I$; warm start $x_0$; inertia range
  $[\underline{w},\overline{w}]=[0.4,0.9]$; cognitive $c_1=1.5$; social $c_2=2.5$
\Ensure $x^\star=\arg\min \mathcal{Q}$ over the visited points
\State draw $x_1,\dots,x_{P-1}$ from a scrambled Sobol net on $[\ell,u]$;\;
  $x_P\gets\Pi_{[\ell,u]}(x_0)$
  \Comment{the warm start is always a particle}
\State $v_i\gets 0$; \; $b_i\gets x_i$, $f_i\gets\mathcal{Q}(x_i)$ evaluated in
  parallel; \; $g\gets\arg\min_i f_i$
\For{$t=1,\dots,I$}
    \State $w\gets\overline{w}-(\overline{w}-\underline{w})\,t/I$
      \Comment{inertia declines linearly}
    \For{$i=1,\dots,P$ \textbf{in parallel}}
        \State draw $r_1,r_2\sim U[0,1]^d$
        \State $v_i\gets w\,v_i+c_1 r_1\odot(b_i-x_i)+c_2 r_2\odot(g-x_i)$
        \State $v_i\gets\Pi_{[-\bar v,\bar v]}(v_i)$, \; $\bar v=0.1\,(u-\ell)$
          \Comment{velocity clamping}
        \State $x_i\gets x_i+v_i$; reflect $x_i$ back into $[\ell,u]$, damping the
          reflected velocity
        \State enforce the parameter constraints (monotone distance bins when
          $\alpha$ is binned)
        \State $f\gets\mathcal{Q}(x_i)$
          \Comment{one call to Algorithm~\ref{alg:simulate}}
        \State \textbf{if} $f<\mathcal{Q}(b_i)$ \textbf{then} $b_i\gets x_i$;\;
          \textbf{if} $f<\mathcal{Q}(g)$ \textbf{then} $g\gets x_i$
    \EndFor
\EndFor
\State \Return $g$
\end{algorithmic}
\end{algorithm}

% ---------------------------------------------------------------------
%  Algorithm 5 -- estimation
% ---------------------------------------------------------------------
\begin{algorithm}[H]
\caption{SMM estimation of $\mu$ by particle swarm optimisation}
\label{alg:smm}
\begin{algorithmic}[1]
\Require empirical moments $\widehat{M}$; empirical sourcing shares
  $\widehat{\gamma}_{as}$ and counts $\widehat{G}_s(0)$; bootstrap covariance
  $\Sigma_{data}$; number of refinement loops $L$; swarm size $P$; iterations $I$
\Ensure $\widehat{\mu}_2$, standard errors, and the calibrated
  $\widehat{T},\widehat{N}$ with their standard errors
\vspace{2pt}
\State \textbf{Starting values.} $\Omega_L,\{\Omega_s\}$ at their input--output
  counterparts; $A_r\propto\pi_r^{1/\varepsilon}$; $\alpha$ at the
  extensive-margin coefficient divided by $\theta$. Call this $\mu^{(0)}$, and
  set the box $[\ell,u]$ multiplicatively around it.
\vspace{3pt}
\State \textbf{Step 1: first fit under equal weights.} Let
  $\mathcal{Q}_I(\mu)=\bigl(\widetilde{M}(\mu)-\widehat{M}\bigr)'\bigl(\widetilde{M}(\mu)-\widehat{M}\bigr)$,
  each evaluation running Algorithm~\ref{alg:simulate}.
\State \quad $\mu\gets\Call{Pso}{\mathcal{Q}_I,[\ell,u],P,I,\mu^{(0)}}$
  \Comment{joint fit over all of $\mu$}
\For{$\text{loop}=1,\dots,L$}
    \Comment{alternating block refinement, search radius shrinking with the loop}
    \For{$\mathcal{B}\in\bigl(\{A_r\},\ \{\alpha\},\ \{\Omega_L,\Omega_s\}\bigr)$}
        \State $\mu\gets\Call{Pso}{\mathcal{Q}_I,\mathcal{B}\text{-box of radius }
          \varrho_{\text{loop}}\text{ around }\mu,P,I,\mu}$
        \Comment{coordinates outside $\mathcal{B}$ fixed}
    \EndFor
    \State \textbf{break if} $\max_k|\Delta\mu_k|/|\mu_k|<\epsilon_\mu$
\EndFor
\State $\widehat{\mu}_1\gets\mu$
\vspace{3pt}
\State \textbf{Step 2: the weighting matrix.} At $\widehat{\mu}_1$, re-simulate
  $K$ times with independent draws, $\Sigma_{sim}\gets\widehat{\mathrm{Var}}
  \bigl(\widetilde{M}^{(1)},\dots,\widetilde{M}^{(K)}\bigr)$, and set
  $W\gets(\Sigma_{data}+\Sigma_{sim})^{-1}$.
\vspace{3pt}
\State \textbf{Step 3: the efficient fit.} Let
  $\mathcal{Q}_W(\mu)=\bigl(\widetilde{M}(\mu)-\widehat{M}\bigr)'W
  \bigl(\widetilde{M}(\mu)-\widehat{M}\bigr)$ and repeat the joint fit and the block
  refinement of Step~1, starting from
  $\widehat{\mu}_1$, yielding $\widehat{\mu}_2$.
\vspace{3pt}
\State \textbf{Step 4: inference.} At $\widehat{\mu}_2$:
\State \quad $G\gets\partial\widetilde{M}/\partial\mu$ by central finite
  differences with common random numbers
\State \quad report $(G'WG)^{-1}$ and
  $(G'WG)^{-1}G'W\Omega WG(G'WG)^{-1}$, $\Omega=\Sigma_{data}+\Sigma_{sim}$
\State \quad $\widehat{T}$: delta method on the calibration map
  $T^\star(\alpha,\widetilde{\gamma})$, equation~(B31)
\State \quad $\widehat{N}$: delta method on
  $\widetilde{G}_s(\widehat{N}_s)=\widehat{G}_s(0)$, equation~(B32)
\State \Return $\widehat{\mu}_2$ and the associated standard errors
\end{algorithmic}
\end{algorithm}

Two implementation notes. The block refinement inside Step~1 exists because the
blocks of $\mu$ have very different curvature: $\{A_r\}$ is pinned almost
one-for-one by the regional sales shares, while $\alpha$ is identified off the
extensive margin, and a single joint swarm spends most of its budget on the
former. Each sub-stage re-runs the swarm on one block, warm-started at the
incumbent and floored at its criterion value, so the criterion is monotone
across sub-stages by construction. The efficient fit of Step~3 restricts the
criterion to the moment blocks that carry the bootstrap covariance --- the
extensive-margin coefficients, the area-aggregated sourcing shares and the
count moments --- and re-optimises $\alpha$ within them, holding
$\{A_r,\Omega_L,\Omega_s\}$ at $\widehat{\mu}_1$; the remaining blocks are the
input--output and sales targets, which are treated as fixed and carry no
sampling variance.
 it where you want the
%  algorithms to appear (e.g. right after Appendix B.3.2).
%
%  Required packages in the master preamble:
%      \usepackage{amsmath,amssymb}
%      \usepackage{algorithm}
%      \usepackage{algpseudocode}
%      \usepackage{float}          % for the [H] float placement used below
%
%  Cross-references: equation numbers are written literally as (B12), (B22),
%  (B31), (B32) to match the current draft. Replace them by \eqref{...} once
%  those equations carry labels.
%  Recommended (renames Require/Ensure to Input/Output):
%      \renewcommand{\algorithmicrequire}{\textbf{Input:}}
%      \renewcommand{\algorithmicensure}{\textbf{Output:}}
%
%  `pseudocode.tex` in the same folder is a standalone driver that compiles
%  this file on its own, for checking.
% =====================================================================

\subsection{Pseudo-code}
\label{app:pseudocode}

This section states the two algorithms of Appendix~B.3 in
explicit form: Algorithm~\ref{alg:simulate} evaluates the model at a candidate
parameter vector $\mu$ and returns the simulated moment vector
$\widetilde{M}(\mu)$, and Algorithm~\ref{alg:smm} minimises the SMM criterion
over $\mu$ by particle swarm optimisation. Algorithms~\ref{alg:invertT},
\ref{alg:countN} and \ref{alg:pso} are the three subroutines they call: the
comparative-advantage inversion, the variety-count calibration, and the swarm
itself.

Throughout, $\mu=\bigl(\Omega_L,\{\Omega_s\},\alpha,\{A_r\}\bigr)$ is the vector
searched over; $a(r')$ denotes the attraction area of region $r'$;
$\mathcal{R}_a$ the regions of area $a$; $\mathcal{R}_D$ the set of downstream
buyers, $R_D=|\mathcal{R}_D|$; and $\mathcal{R}^+_s$ the regions inside active
attraction areas for sector $s$. The calibrated objects
\[
  \bigl\{\theta,\ \varepsilon,\ \lambda,\ \nu,\ \{\nu_s\},\
  \{w_{r's}\},\ \{\delta_r\},\ \{d_{r'r}\}\bigr\}
\]
are fixed once and for all and are suppressed from the argument lists.

% ---------------------------------------------------------------------
%  Algorithm 1 -- simulation
% ---------------------------------------------------------------------
\begin{algorithm}[H]
\caption{Simulation of the economy at a candidate $\mu$ (one criterion evaluation)}
\label{alg:simulate}
\begin{algorithmic}[1]
\Require candidate $\mu=(\Omega_L,\{\Omega_s\},\alpha,\{A_r\})$; empirical
  sourcing shares $\widehat{\gamma}_{as}$; empirical count moments
  $\widehat{G}_s(0)$ and bounds $[N^{lo}_s,N^{hi}_s]$; a fixed array of uniform
  draws $u_{r'\rho s}\in(0,1)$, $\rho=1,\dots,m$
\Ensure simulated moment vector $\widetilde{M}(\mu)$, and the calibrated
  $\{T_{as}\},\{\widehat{N}_s\}$
\vspace{2pt}
\State $\tau_{r'r}\gets d_{r'r}^{\,\alpha}$
  \Comment{trade costs; the only place $\alpha$ enters}
\vspace{2pt}
\State \textbf{Step 1 (comparative advantage).}
  $\{T_{as}\}\gets\Call{InvertT}{\alpha,\Omega_L,\{\Omega_s\},\{A_r\};\widehat{\gamma}}$
  \Comment{Algorithm~\ref{alg:invertT}}
\vspace{2pt}
\State \textbf{Step 2 (Ricardian competition).}
\For{$s=1,\dots,S$ \textbf{and} $\rho=1,\dots,m$}
    \For{$r'=1,\dots,R$}
        \State $z_{r'\rho s}\gets T_{a(r')s}^{1/\theta}\bigl(-\log u_{r'\rho s}\bigr)^{-1/\theta}$
          \Comment{Fr\'echet inverse c.d.f.}
    \EndFor
    \For{$r\in\mathcal{R}_D$}
        \State $r^\star(\rho,s,r)\gets\arg\min_{r'}\; w_{r's}\tau_{r'r}/z_{r'\rho s}$
          \Comment{winner of variety $\rho$ in market $(r,s)$}
        \State $p_{\rho s r}\gets w_{r^\star s}\tau_{r^\star r}/z_{r^\star\rho s}$
    \EndFor
\EndFor
\State $k_{r's}\gets\#\{\rho:\ r'=r^\star(\rho,s,r)\ \text{for some }r\in\mathcal{R}_D\}$,
  \quad $p_{r's}\gets k_{r's}/m$
  \Comment{win-anywhere counts and probabilities}
\vspace{2pt}
\State \textbf{Prices, sales and flows.} For every $r\in\mathcal{R}_D$:
\State \quad $P_{rs}\gets\bigl[\tfrac1m\sum_{\rho}p_{\rho s r}^{\,1-\nu_s}\bigr]^{1/(1-\nu_s)}$,\quad
  $P_r\gets\bigl[\sum_s\Omega_s P_{rs}^{1-\nu}\bigr]^{1/(1-\nu)}$
\State \quad $c_r\gets\bigl[\Omega_L w_r^{1-\lambda}+(1-\Omega_L)P_r^{1-\lambda}\bigr]^{1/(1-\lambda)}$,\quad
  $\tilde c_r\gets c_r/A_r$
\State \quad $y_r\gets \tilde c_r^{\,\varepsilon}\,\delta_r\big/\sum_{r''}\tilde c_{r''}^{\,\varepsilon}\delta_{r''}$,\quad
  $\pi_r\gets y_r$
  \Comment{iso-elastic demand, $E$ normalised to $1$}
\State \quad $X_{r'rs}\gets$ expenditure of $r$ on the varieties won by $r'$,
  from~(B12)
\vspace{2pt}
\State \textbf{Step 3 (variety count).}
  $\{\widehat{N}_s\}\gets\Call{CalibrateN}{\{k_{r's}\},m;\widehat{G}_s(0),[N^{lo}_s,N^{hi}_s]}$
  \Comment{Algorithm~\ref{alg:countN}}
\vspace{2pt}
\State \textbf{Step 4 (moments).} Assemble $\widetilde{M}(\mu)$ from the single
  solve above:
\State \quad (1) $\widetilde{\Omega}_L$, (2) $\widetilde{\Omega}_s$, (3) $\pi_r$
  \Comment{value blocks, from $\{X_{r'rs}\},\{c_r\},\{y_r\}$}
\State \quad (4) $\beta$: cloglog of ``does not supply'' on $\log d_{r'a(r')}$ and
  $\log z$, area${}\times{}$sector fixed effects, run on the simulated firm
  population $\{(r',\rho,s)\}$
\State \quad (5) $\widetilde{\gamma}_{as}\gets\sum_{r'\in\mathcal{R}_a}\sum_r X_{r'rs}\big/\sum_{r''}\sum_r X_{r''rs}$
\State \quad (6) $\widetilde{G}_s(\widehat{N}_s)$
  \Comment{count moment, evaluated at the calibrated $\widehat{N}_s$}
\State \Return $\widetilde{M}(\mu)$
\end{algorithmic}
\end{algorithm}

Three features of Algorithm~\ref{alg:simulate} are worth emphasising. First, the
arrows of~(B22) run one way: $T$ is calibrated without knowing
$N$, and $N$ does not revise $T$, so Steps~1--3 form a recursion rather
than a simultaneous system. Second, the simulation size $m$ is a Monte Carlo
device and is unrelated to $\widehat{N}_s$: the winner of a variety depends only
on relative production costs, so $p_{r's}$ can be estimated at any $m$ large
enough. Third, the uniform draws $u$ are held fixed across every evaluation of
the criterion within an optimisation stage (common random numbers), which is
what makes the criterion a deterministic function of $\mu$ that a swarm can
descend; they are scrambled Sobol sequences, constructed one low-dimensional
sequence per sector with regions as dimensions and varieties as points.
Independent pseudo-random draws are used instead when computing $\Sigma_{sim}$
and the numerical derivatives, where independence across replications is
required.

% ---------------------------------------------------------------------
%  Algorithm 2 -- comparative advantage inversion
% ---------------------------------------------------------------------
\begin{algorithm}[H]
\caption{\textsc{InvertT}: comparative advantage by matrix scaling with
general-equilibrium expenditure weights}
\label{alg:invertT}
\begin{algorithmic}[1]
\Require $\alpha,\Omega_L,\{\Omega_s\},\{A_r\}$; targets $\widehat{\gamma}_{as}$;
  damping $\varsigma\in(0,1]$; tolerance $\epsilon_T$; starting values $T^{(0)}_{as}$
\Ensure $\{T_{as}\}$ reproducing $\widehat{\gamma}_{as}$ exactly
\State $T\gets T^{(0)}$; normalise $T_{as}\gets T_{as}/T_{a_s^{\text{ref}}s}$ for each $s$
  \Comment{one reference area per sector fixes the scale}
\For{$t=1,2,\dots$ \textbf{until} convergence}
    \State $B_{ar}\gets\sum_{r'\in\mathcal{R}_a}(w_{r's}\tau_{r'r})^{-\theta}$,\quad
      $\Phi_{rs}\gets\sum_{a'}T_{a's}B_{a'r}$
    \State $\omega_{rs}\gets$ expenditure weights implied by $T$
      \Comment{price and sales block of Algorithm~\ref{alg:simulate}}
    \State $\gamma_{as}(T)\gets\sum_r\omega_{rs}\,T_{as}B_{ar}/\Phi_{rs}$
    \State $T^{+}_{as}\gets T_{as}\bigl(\widehat{\gamma}_{as}/\gamma_{as}(T)\bigr)^{\varsigma}$
      \Comment{multiplicative Sinkhorn update, in logs}
    \State $T^{+}_{as}\gets T^{+}_{as}/T^{+}_{a^{\text{ref}}_s s}$ for each $s$
    \State $\text{resid}\gets\max_{a,s}\bigl|\log T^{+}_{as}-\log T_{as}\bigr|$;\quad $T\gets T^{+}$
    \State \textbf{if} $\text{resid}<\epsilon_T$ \textbf{then break}
\EndFor
\State \Return $T$
\end{algorithmic}
\end{algorithm}

With expenditure weights held fixed, the multiplicative update is a Sinkhorn scaling of the strictly positive matrix
$B$, hence a contraction with a fixed point unique up to the sector-specific
scale removed by the reference-area normalisation. Recomputing $\omega_{rs}$
inside the loop adds the general-equilibrium response of prices and sales, so
the full map is a fixed-point iteration; the damping $\varsigma$ ($0.9$ in the
implementation) and the stopping rule on $\max_{a,s}|\Delta\log T_{as}|$ are what
the text refers to as requiring both sourcing shares and comparative advantages
to converge.

% ---------------------------------------------------------------------
%  Algorithm 3 -- variety count
% ---------------------------------------------------------------------
\begin{algorithm}[H]
\caption{\textsc{CalibrateN}: variety count by monotone integer search}
\label{alg:countN}
\begin{algorithmic}[1]
\Require win counts $k_{r's}$ from $m$ simulated varieties; targets
  $\widehat{G}_s(0)$; bounds $[N^{lo}_s,N^{hi}_s]$
\Ensure $\{\widehat{N}_s\}$
\For{$s=1,\dots,S$}
    \State define
      $\widetilde{G}_s(n)\gets\dfrac{1}{|\mathcal{R}^+_s|}\displaystyle\sum_{r'\in\mathcal{R}^+_s}
      \binom{m-k_{r's}}{n}\Big/\binom{m}{n}$
      \Comment{unbiased; decreasing in $n$}
    \If{$\widetilde{G}_s(N^{lo}_s)\le\widehat{G}_s(0)$}
        \State $\widehat{N}_s\gets N^{lo}_s$ \Comment{clamped at the lower bound}
    \ElsIf{$\widetilde{G}_s(N^{hi}_s)\ge\widehat{G}_s(0)$}
        \State $\widehat{N}_s\gets N^{hi}_s$ \Comment{clamped at the upper bound}
    \Else
        \State $(\underline{n},\overline{n})\gets(N^{lo}_s,N^{hi}_s)$
        \While{$\overline{n}-\underline{n}>1$}
            \State $n\gets\lfloor(\underline{n}+\overline{n})/2\rfloor$
            \State \textbf{if} $\widetilde{G}_s(n)>\widehat{G}_s(0)$ \textbf{then}
              $\underline{n}\gets n$ \textbf{else} $\overline{n}\gets n$
        \EndWhile
        \State $\widehat{N}_s\gets\arg\min_{n\in\{\underline{n},\overline{n}\}}
          \bigl|\widetilde{G}_s(n)-\widehat{G}_s(0)\bigr|$
    \EndIf
\EndFor
\State \Return $\{\widehat{N}_s\}$
\end{algorithmic}
\end{algorithm}

Because $\widetilde{G}_s$ is decreasing in $n$ and is evaluated in closed form
from the winning matrix of Step~2, the search costs $O(\log(N^{hi}_s-N^{lo}_s))$
evaluations and never re-simulates the Ricardian competition. The count is an
integer at every evaluation: it is never relaxed to the real line and never
rounded. A bound that binds is recorded and reported rather than silently
absorbed.

% ---------------------------------------------------------------------
%  Algorithm 4 -- PSO
% ---------------------------------------------------------------------
\begin{algorithm}[H]
\caption{\textsc{Pso}: particle swarm minimisation of $\mathcal{Q}$ over a box}
\label{alg:pso}
\begin{algorithmic}[1]
\Require objective $\mathcal{Q}$; box $[\ell,u]\subset\mathbb{R}^{d}$; swarm size
  $P$; iterations $I$; warm start $x_0$; inertia range
  $[\underline{w},\overline{w}]=[0.4,0.9]$; cognitive $c_1=1.5$; social $c_2=2.5$
\Ensure $x^\star=\arg\min \mathcal{Q}$ over the visited points
\State draw $x_1,\dots,x_{P-1}$ from a scrambled Sobol net on $[\ell,u]$;\;
  $x_P\gets\Pi_{[\ell,u]}(x_0)$
  \Comment{the warm start is always a particle}
\State $v_i\gets 0$; \; $b_i\gets x_i$, $f_i\gets\mathcal{Q}(x_i)$ evaluated in
  parallel; \; $g\gets\arg\min_i f_i$
\For{$t=1,\dots,I$}
    \State $w\gets\overline{w}-(\overline{w}-\underline{w})\,t/I$
      \Comment{inertia declines linearly}
    \For{$i=1,\dots,P$ \textbf{in parallel}}
        \State draw $r_1,r_2\sim U[0,1]^d$
        \State $v_i\gets w\,v_i+c_1 r_1\odot(b_i-x_i)+c_2 r_2\odot(g-x_i)$
        \State $v_i\gets\Pi_{[-\bar v,\bar v]}(v_i)$, \; $\bar v=0.1\,(u-\ell)$
          \Comment{velocity clamping}
        \State $x_i\gets x_i+v_i$; reflect $x_i$ back into $[\ell,u]$, damping the
          reflected velocity
        \State enforce the parameter constraints (monotone distance bins when
          $\alpha$ is binned)
        \State $f\gets\mathcal{Q}(x_i)$
          \Comment{one call to Algorithm~\ref{alg:simulate}}
        \State \textbf{if} $f<\mathcal{Q}(b_i)$ \textbf{then} $b_i\gets x_i$;\;
          \textbf{if} $f<\mathcal{Q}(g)$ \textbf{then} $g\gets x_i$
    \EndFor
\EndFor
\State \Return $g$
\end{algorithmic}
\end{algorithm}

% ---------------------------------------------------------------------
%  Algorithm 5 -- estimation
% ---------------------------------------------------------------------
\begin{algorithm}[H]
\caption{SMM estimation of $\mu$ by particle swarm optimisation}
\label{alg:smm}
\begin{algorithmic}[1]
\Require empirical moments $\widehat{M}$; empirical sourcing shares
  $\widehat{\gamma}_{as}$ and counts $\widehat{G}_s(0)$; bootstrap covariance
  $\Sigma_{data}$; number of refinement loops $L$; swarm size $P$; iterations $I$
\Ensure $\widehat{\mu}_2$, standard errors, and the calibrated
  $\widehat{T},\widehat{N}$ with their standard errors
\vspace{2pt}
\State \textbf{Starting values.} $\Omega_L,\{\Omega_s\}$ at their input--output
  counterparts; $A_r\propto\pi_r^{1/\varepsilon}$; $\alpha$ at the
  extensive-margin coefficient divided by $\theta$. Call this $\mu^{(0)}$, and
  set the box $[\ell,u]$ multiplicatively around it.
\vspace{3pt}
\State \textbf{Step 1: first fit under equal weights.} Let
  $\mathcal{Q}_I(\mu)=\bigl(\widetilde{M}(\mu)-\widehat{M}\bigr)'\bigl(\widetilde{M}(\mu)-\widehat{M}\bigr)$,
  each evaluation running Algorithm~\ref{alg:simulate}.
\State \quad $\mu\gets\Call{Pso}{\mathcal{Q}_I,[\ell,u],P,I,\mu^{(0)}}$
  \Comment{joint fit over all of $\mu$}
\For{$\text{loop}=1,\dots,L$}
    \Comment{alternating block refinement, search radius shrinking with the loop}
    \For{$\mathcal{B}\in\bigl(\{A_r\},\ \{\alpha\},\ \{\Omega_L,\Omega_s\}\bigr)$}
        \State $\mu\gets\Call{Pso}{\mathcal{Q}_I,\mathcal{B}\text{-box of radius }
          \varrho_{\text{loop}}\text{ around }\mu,P,I,\mu}$
        \Comment{coordinates outside $\mathcal{B}$ fixed}
    \EndFor
    \State \textbf{break if} $\max_k|\Delta\mu_k|/|\mu_k|<\epsilon_\mu$
\EndFor
\State $\widehat{\mu}_1\gets\mu$
\vspace{3pt}
\State \textbf{Step 2: the weighting matrix.} At $\widehat{\mu}_1$, re-simulate
  $K$ times with independent draws, $\Sigma_{sim}\gets\widehat{\mathrm{Var}}
  \bigl(\widetilde{M}^{(1)},\dots,\widetilde{M}^{(K)}\bigr)$, and set
  $W\gets(\Sigma_{data}+\Sigma_{sim})^{-1}$.
\vspace{3pt}
\State \textbf{Step 3: the efficient fit.} Let
  $\mathcal{Q}_W(\mu)=\bigl(\widetilde{M}(\mu)-\widehat{M}\bigr)'W
  \bigl(\widetilde{M}(\mu)-\widehat{M}\bigr)$ and repeat the joint fit and the block
  refinement of Step~1, starting from
  $\widehat{\mu}_1$, yielding $\widehat{\mu}_2$.
\vspace{3pt}
\State \textbf{Step 4: inference.} At $\widehat{\mu}_2$:
\State \quad $G\gets\partial\widetilde{M}/\partial\mu$ by central finite
  differences with common random numbers
\State \quad report $(G'WG)^{-1}$ and
  $(G'WG)^{-1}G'W\Omega WG(G'WG)^{-1}$, $\Omega=\Sigma_{data}+\Sigma_{sim}$
\State \quad $\widehat{T}$: delta method on the calibration map
  $T^\star(\alpha,\widetilde{\gamma})$, equation~(B31)
\State \quad $\widehat{N}$: delta method on
  $\widetilde{G}_s(\widehat{N}_s)=\widehat{G}_s(0)$, equation~(B32)
\State \Return $\widehat{\mu}_2$ and the associated standard errors
\end{algorithmic}
\end{algorithm}

Two implementation notes. The block refinement inside Step~1 exists because the
blocks of $\mu$ have very different curvature: $\{A_r\}$ is pinned almost
one-for-one by the regional sales shares, while $\alpha$ is identified off the
extensive margin, and a single joint swarm spends most of its budget on the
former. Each sub-stage re-runs the swarm on one block, warm-started at the
incumbent and floored at its criterion value, so the criterion is monotone
across sub-stages by construction. The efficient fit of Step~3 restricts the
criterion to the moment blocks that carry the bootstrap covariance --- the
extensive-margin coefficients, the area-aggregated sourcing shares and the
count moments --- and re-optimises $\alpha$ within them, holding
$\{A_r,\Omega_L,\Omega_s\}$ at $\widehat{\mu}_1$; the remaining blocks are the
input--output and sales targets, which are treated as fixed and carry no
sampling variance.
 it where you want the
%  algorithms to appear (e.g. right after Appendix B.3.2).
%
%  Required packages in the master preamble:
%      \usepackage{amsmath,amssymb}
%      \usepackage{algorithm}
%      \usepackage{algpseudocode}
%      \usepackage{float}          % for the [H] float placement used below
%
%  Cross-references: equation numbers are written literally as (B12), (B22),
%  (B31), (B32) to match the current draft. Replace them by \eqref{...} once
%  those equations carry labels.
%  Recommended (renames Require/Ensure to Input/Output):
%      \renewcommand{\algorithmicrequire}{\textbf{Input:}}
%      \renewcommand{\algorithmicensure}{\textbf{Output:}}
%
%  `pseudocode.tex` in the same folder is a standalone driver that compiles
%  this file on its own, for checking.
% =====================================================================

\subsection{Pseudo-code}
\label{app:pseudocode}

This section states the two algorithms of Appendix~B.3 in
explicit form: Algorithm~\ref{alg:simulate} evaluates the model at a candidate
parameter vector $\mu$ and returns the simulated moment vector
$\widetilde{M}(\mu)$, and Algorithm~\ref{alg:smm} minimises the SMM criterion
over $\mu$ by particle swarm optimisation. Algorithms~\ref{alg:invertT},
\ref{alg:countN} and \ref{alg:pso} are the three subroutines they call: the
comparative-advantage inversion, the variety-count calibration, and the swarm
itself.

Throughout, $\mu=\bigl(\Omega_L,\{\Omega_s\},\alpha,\{A_r\}\bigr)$ is the vector
searched over; $a(r')$ denotes the attraction area of region $r'$;
$\mathcal{R}_a$ the regions of area $a$; $\mathcal{R}_D$ the set of downstream
buyers, $R_D=|\mathcal{R}_D|$; and $\mathcal{R}^+_s$ the regions inside active
attraction areas for sector $s$. The calibrated objects
\[
  \bigl\{\theta,\ \varepsilon,\ \lambda,\ \nu,\ \{\nu_s\},\
  \{w_{r's}\},\ \{\delta_r\},\ \{d_{r'r}\}\bigr\}
\]
are fixed once and for all and are suppressed from the argument lists.

% ---------------------------------------------------------------------
%  Algorithm 1 -- simulation
% ---------------------------------------------------------------------
\begin{algorithm}[H]
\caption{Simulation of the economy at a candidate $\mu$ (one criterion evaluation)}
\label{alg:simulate}
\begin{algorithmic}[1]
\Require candidate $\mu=(\Omega_L,\{\Omega_s\},\alpha,\{A_r\})$; empirical
  sourcing shares $\widehat{\gamma}_{as}$; empirical count moments
  $\widehat{G}_s(0)$ and bounds $[N^{lo}_s,N^{hi}_s]$; a fixed array of uniform
  draws $u_{r'\rho s}\in(0,1)$, $\rho=1,\dots,m$
\Ensure simulated moment vector $\widetilde{M}(\mu)$, and the calibrated
  $\{T_{as}\},\{\widehat{N}_s\}$
\vspace{2pt}
\State $\tau_{r'r}\gets d_{r'r}^{\,\alpha}$
  \Comment{trade costs; the only place $\alpha$ enters}
\vspace{2pt}
\State \textbf{Step 1 (comparative advantage).}
  $\{T_{as}\}\gets\Call{InvertT}{\alpha,\Omega_L,\{\Omega_s\},\{A_r\};\widehat{\gamma}}$
  \Comment{Algorithm~\ref{alg:invertT}}
\vspace{2pt}
\State \textbf{Step 2 (Ricardian competition).}
\For{$s=1,\dots,S$ \textbf{and} $\rho=1,\dots,m$}
    \For{$r'=1,\dots,R$}
        \State $z_{r'\rho s}\gets T_{a(r')s}^{1/\theta}\bigl(-\log u_{r'\rho s}\bigr)^{-1/\theta}$
          \Comment{Fr\'echet inverse c.d.f.}
    \EndFor
    \For{$r\in\mathcal{R}_D$}
        \State $r^\star(\rho,s,r)\gets\arg\min_{r'}\; w_{r's}\tau_{r'r}/z_{r'\rho s}$
          \Comment{winner of variety $\rho$ in market $(r,s)$}
        \State $p_{\rho s r}\gets w_{r^\star s}\tau_{r^\star r}/z_{r^\star\rho s}$
    \EndFor
\EndFor
\State $k_{r's}\gets\#\{\rho:\ r'=r^\star(\rho,s,r)\ \text{for some }r\in\mathcal{R}_D\}$,
  \quad $p_{r's}\gets k_{r's}/m$
  \Comment{win-anywhere counts and probabilities}
\vspace{2pt}
\State \textbf{Prices, sales and flows.} For every $r\in\mathcal{R}_D$:
\State \quad $P_{rs}\gets\bigl[\tfrac1m\sum_{\rho}p_{\rho s r}^{\,1-\nu_s}\bigr]^{1/(1-\nu_s)}$,\quad
  $P_r\gets\bigl[\sum_s\Omega_s P_{rs}^{1-\nu}\bigr]^{1/(1-\nu)}$
\State \quad $c_r\gets\bigl[\Omega_L w_r^{1-\lambda}+(1-\Omega_L)P_r^{1-\lambda}\bigr]^{1/(1-\lambda)}$,\quad
  $\tilde c_r\gets c_r/A_r$
\State \quad $y_r\gets \tilde c_r^{\,\varepsilon}\,\delta_r\big/\sum_{r''}\tilde c_{r''}^{\,\varepsilon}\delta_{r''}$,\quad
  $\pi_r\gets y_r$
  \Comment{iso-elastic demand, $E$ normalised to $1$}
\State \quad $X_{r'rs}\gets$ expenditure of $r$ on the varieties won by $r'$,
  from~(B12)
\vspace{2pt}
\State \textbf{Step 3 (variety count).}
  $\{\widehat{N}_s\}\gets\Call{CalibrateN}{\{k_{r's}\},m;\widehat{G}_s(0),[N^{lo}_s,N^{hi}_s]}$
  \Comment{Algorithm~\ref{alg:countN}}
\vspace{2pt}
\State \textbf{Step 4 (moments).} Assemble $\widetilde{M}(\mu)$ from the single
  solve above:
\State \quad (1) $\widetilde{\Omega}_L$, (2) $\widetilde{\Omega}_s$, (3) $\pi_r$
  \Comment{value blocks, from $\{X_{r'rs}\},\{c_r\},\{y_r\}$}
\State \quad (4) $\beta$: cloglog of ``does not supply'' on $\log d_{r'a(r')}$ and
  $\log z$, area${}\times{}$sector fixed effects, run on the simulated firm
  population $\{(r',\rho,s)\}$
\State \quad (5) $\widetilde{\gamma}_{as}\gets\sum_{r'\in\mathcal{R}_a}\sum_r X_{r'rs}\big/\sum_{r''}\sum_r X_{r''rs}$
\State \quad (6) $\widetilde{G}_s(\widehat{N}_s)$
  \Comment{count moment, evaluated at the calibrated $\widehat{N}_s$}
\State \Return $\widetilde{M}(\mu)$
\end{algorithmic}
\end{algorithm}

Three features of Algorithm~\ref{alg:simulate} are worth emphasising. First, the
arrows of~(B22) run one way: $T$ is calibrated without knowing
$N$, and $N$ does not revise $T$, so Steps~1--3 form a recursion rather
than a simultaneous system. Second, the simulation size $m$ is a Monte Carlo
device and is unrelated to $\widehat{N}_s$: the winner of a variety depends only
on relative production costs, so $p_{r's}$ can be estimated at any $m$ large
enough. Third, the uniform draws $u$ are held fixed across every evaluation of
the criterion within an optimisation stage (common random numbers), which is
what makes the criterion a deterministic function of $\mu$ that a swarm can
descend; they are scrambled Sobol sequences, constructed one low-dimensional
sequence per sector with regions as dimensions and varieties as points.
Independent pseudo-random draws are used instead when computing $\Sigma_{sim}$
and the numerical derivatives, where independence across replications is
required.

% ---------------------------------------------------------------------
%  Algorithm 2 -- comparative advantage inversion
% ---------------------------------------------------------------------
\begin{algorithm}[H]
\caption{\textsc{InvertT}: comparative advantage by matrix scaling with
general-equilibrium expenditure weights}
\label{alg:invertT}
\begin{algorithmic}[1]
\Require $\alpha,\Omega_L,\{\Omega_s\},\{A_r\}$; targets $\widehat{\gamma}_{as}$;
  damping $\varsigma\in(0,1]$; tolerance $\epsilon_T$; starting values $T^{(0)}_{as}$
\Ensure $\{T_{as}\}$ reproducing $\widehat{\gamma}_{as}$ exactly
\State $T\gets T^{(0)}$; normalise $T_{as}\gets T_{as}/T_{a_s^{\text{ref}}s}$ for each $s$
  \Comment{one reference area per sector fixes the scale}
\For{$t=1,2,\dots$ \textbf{until} convergence}
    \State $B_{ar}\gets\sum_{r'\in\mathcal{R}_a}(w_{r's}\tau_{r'r})^{-\theta}$,\quad
      $\Phi_{rs}\gets\sum_{a'}T_{a's}B_{a'r}$
    \State $\omega_{rs}\gets$ expenditure weights implied by $T$
      \Comment{price and sales block of Algorithm~\ref{alg:simulate}}
    \State $\gamma_{as}(T)\gets\sum_r\omega_{rs}\,T_{as}B_{ar}/\Phi_{rs}$
    \State $T^{+}_{as}\gets T_{as}\bigl(\widehat{\gamma}_{as}/\gamma_{as}(T)\bigr)^{\varsigma}$
      \Comment{multiplicative Sinkhorn update, in logs}
    \State $T^{+}_{as}\gets T^{+}_{as}/T^{+}_{a^{\text{ref}}_s s}$ for each $s$
    \State $\text{resid}\gets\max_{a,s}\bigl|\log T^{+}_{as}-\log T_{as}\bigr|$;\quad $T\gets T^{+}$
    \State \textbf{if} $\text{resid}<\epsilon_T$ \textbf{then break}
\EndFor
\State \Return $T$
\end{algorithmic}
\end{algorithm}

With expenditure weights held fixed, the multiplicative update is a Sinkhorn scaling of the strictly positive matrix
$B$, hence a contraction with a fixed point unique up to the sector-specific
scale removed by the reference-area normalisation. Recomputing $\omega_{rs}$
inside the loop adds the general-equilibrium response of prices and sales, so
the full map is a fixed-point iteration; the damping $\varsigma$ ($0.9$ in the
implementation) and the stopping rule on $\max_{a,s}|\Delta\log T_{as}|$ are what
the text refers to as requiring both sourcing shares and comparative advantages
to converge.

% ---------------------------------------------------------------------
%  Algorithm 3 -- variety count
% ---------------------------------------------------------------------
\begin{algorithm}[H]
\caption{\textsc{CalibrateN}: variety count by monotone integer search}
\label{alg:countN}
\begin{algorithmic}[1]
\Require win counts $k_{r's}$ from $m$ simulated varieties; targets
  $\widehat{G}_s(0)$; bounds $[N^{lo}_s,N^{hi}_s]$
\Ensure $\{\widehat{N}_s\}$
\For{$s=1,\dots,S$}
    \State define
      $\widetilde{G}_s(n)\gets\dfrac{1}{|\mathcal{R}^+_s|}\displaystyle\sum_{r'\in\mathcal{R}^+_s}
      \binom{m-k_{r's}}{n}\Big/\binom{m}{n}$
      \Comment{unbiased; decreasing in $n$}
    \If{$\widetilde{G}_s(N^{lo}_s)\le\widehat{G}_s(0)$}
        \State $\widehat{N}_s\gets N^{lo}_s$ \Comment{clamped at the lower bound}
    \ElsIf{$\widetilde{G}_s(N^{hi}_s)\ge\widehat{G}_s(0)$}
        \State $\widehat{N}_s\gets N^{hi}_s$ \Comment{clamped at the upper bound}
    \Else
        \State $(\underline{n},\overline{n})\gets(N^{lo}_s,N^{hi}_s)$
        \While{$\overline{n}-\underline{n}>1$}
            \State $n\gets\lfloor(\underline{n}+\overline{n})/2\rfloor$
            \State \textbf{if} $\widetilde{G}_s(n)>\widehat{G}_s(0)$ \textbf{then}
              $\underline{n}\gets n$ \textbf{else} $\overline{n}\gets n$
        \EndWhile
        \State $\widehat{N}_s\gets\arg\min_{n\in\{\underline{n},\overline{n}\}}
          \bigl|\widetilde{G}_s(n)-\widehat{G}_s(0)\bigr|$
    \EndIf
\EndFor
\State \Return $\{\widehat{N}_s\}$
\end{algorithmic}
\end{algorithm}

Because $\widetilde{G}_s$ is decreasing in $n$ and is evaluated in closed form
from the winning matrix of Step~2, the search costs $O(\log(N^{hi}_s-N^{lo}_s))$
evaluations and never re-simulates the Ricardian competition. The count is an
integer at every evaluation: it is never relaxed to the real line and never
rounded. A bound that binds is recorded and reported rather than silently
absorbed.

% ---------------------------------------------------------------------
%  Algorithm 4 -- PSO
% ---------------------------------------------------------------------
\begin{algorithm}[H]
\caption{\textsc{Pso}: particle swarm minimisation of $\mathcal{Q}$ over a box}
\label{alg:pso}
\begin{algorithmic}[1]
\Require objective $\mathcal{Q}$; box $[\ell,u]\subset\mathbb{R}^{d}$; swarm size
  $P$; iterations $I$; warm start $x_0$; inertia range
  $[\underline{w},\overline{w}]=[0.4,0.9]$; cognitive $c_1=1.5$; social $c_2=2.5$
\Ensure $x^\star=\arg\min \mathcal{Q}$ over the visited points
\State draw $x_1,\dots,x_{P-1}$ from a scrambled Sobol net on $[\ell,u]$;\;
  $x_P\gets\Pi_{[\ell,u]}(x_0)$
  \Comment{the warm start is always a particle}
\State $v_i\gets 0$; \; $b_i\gets x_i$, $f_i\gets\mathcal{Q}(x_i)$ evaluated in
  parallel; \; $g\gets\arg\min_i f_i$
\For{$t=1,\dots,I$}
    \State $w\gets\overline{w}-(\overline{w}-\underline{w})\,t/I$
      \Comment{inertia declines linearly}
    \For{$i=1,\dots,P$ \textbf{in parallel}}
        \State draw $r_1,r_2\sim U[0,1]^d$
        \State $v_i\gets w\,v_i+c_1 r_1\odot(b_i-x_i)+c_2 r_2\odot(g-x_i)$
        \State $v_i\gets\Pi_{[-\bar v,\bar v]}(v_i)$, \; $\bar v=0.1\,(u-\ell)$
          \Comment{velocity clamping}
        \State $x_i\gets x_i+v_i$; reflect $x_i$ back into $[\ell,u]$, damping the
          reflected velocity
        \State enforce the parameter constraints (monotone distance bins when
          $\alpha$ is binned)
        \State $f\gets\mathcal{Q}(x_i)$
          \Comment{one call to Algorithm~\ref{alg:simulate}}
        \State \textbf{if} $f<\mathcal{Q}(b_i)$ \textbf{then} $b_i\gets x_i$;\;
          \textbf{if} $f<\mathcal{Q}(g)$ \textbf{then} $g\gets x_i$
    \EndFor
\EndFor
\State \Return $g$
\end{algorithmic}
\end{algorithm}

% ---------------------------------------------------------------------
%  Algorithm 5 -- estimation
% ---------------------------------------------------------------------
\begin{algorithm}[H]
\caption{SMM estimation of $\mu$ by particle swarm optimisation}
\label{alg:smm}
\begin{algorithmic}[1]
\Require empirical moments $\widehat{M}$; empirical sourcing shares
  $\widehat{\gamma}_{as}$ and counts $\widehat{G}_s(0)$; bootstrap covariance
  $\Sigma_{data}$; number of refinement loops $L$; swarm size $P$; iterations $I$
\Ensure $\widehat{\mu}_2$, standard errors, and the calibrated
  $\widehat{T},\widehat{N}$ with their standard errors
\vspace{2pt}
\State \textbf{Starting values.} $\Omega_L,\{\Omega_s\}$ at their input--output
  counterparts; $A_r\propto\pi_r^{1/\varepsilon}$; $\alpha$ at the
  extensive-margin coefficient divided by $\theta$. Call this $\mu^{(0)}$, and
  set the box $[\ell,u]$ multiplicatively around it.
\vspace{3pt}
\State \textbf{Step 1: first fit under equal weights.} Let
  $\mathcal{Q}_I(\mu)=\bigl(\widetilde{M}(\mu)-\widehat{M}\bigr)'\bigl(\widetilde{M}(\mu)-\widehat{M}\bigr)$,
  each evaluation running Algorithm~\ref{alg:simulate}.
\State \quad $\mu\gets\Call{Pso}{\mathcal{Q}_I,[\ell,u],P,I,\mu^{(0)}}$
  \Comment{joint fit over all of $\mu$}
\For{$\text{loop}=1,\dots,L$}
    \Comment{alternating block refinement, search radius shrinking with the loop}
    \For{$\mathcal{B}\in\bigl(\{A_r\},\ \{\alpha\},\ \{\Omega_L,\Omega_s\}\bigr)$}
        \State $\mu\gets\Call{Pso}{\mathcal{Q}_I,\mathcal{B}\text{-box of radius }
          \varrho_{\text{loop}}\text{ around }\mu,P,I,\mu}$
        \Comment{coordinates outside $\mathcal{B}$ fixed}
    \EndFor
    \State \textbf{break if} $\max_k|\Delta\mu_k|/|\mu_k|<\epsilon_\mu$
\EndFor
\State $\widehat{\mu}_1\gets\mu$
\vspace{3pt}
\State \textbf{Step 2: the weighting matrix.} At $\widehat{\mu}_1$, re-simulate
  $K$ times with independent draws, $\Sigma_{sim}\gets\widehat{\mathrm{Var}}
  \bigl(\widetilde{M}^{(1)},\dots,\widetilde{M}^{(K)}\bigr)$, and set
  $W\gets(\Sigma_{data}+\Sigma_{sim})^{-1}$.
\vspace{3pt}
\State \textbf{Step 3: the efficient fit.} Let
  $\mathcal{Q}_W(\mu)=\bigl(\widetilde{M}(\mu)-\widehat{M}\bigr)'W
  \bigl(\widetilde{M}(\mu)-\widehat{M}\bigr)$ and repeat the joint fit and the block
  refinement of Step~1, starting from
  $\widehat{\mu}_1$, yielding $\widehat{\mu}_2$.
\vspace{3pt}
\State \textbf{Step 4: inference.} At $\widehat{\mu}_2$:
\State \quad $G\gets\partial\widetilde{M}/\partial\mu$ by central finite
  differences with common random numbers
\State \quad report $(G'WG)^{-1}$ and
  $(G'WG)^{-1}G'W\Omega WG(G'WG)^{-1}$, $\Omega=\Sigma_{data}+\Sigma_{sim}$
\State \quad $\widehat{T}$: delta method on the calibration map
  $T^\star(\alpha,\widetilde{\gamma})$, equation~(B31)
\State \quad $\widehat{N}$: delta method on
  $\widetilde{G}_s(\widehat{N}_s)=\widehat{G}_s(0)$, equation~(B32)
\State \Return $\widehat{\mu}_2$ and the associated standard errors
\end{algorithmic}
\end{algorithm}

Two implementation notes. The block refinement inside Step~1 exists because the
blocks of $\mu$ have very different curvature: $\{A_r\}$ is pinned almost
one-for-one by the regional sales shares, while $\alpha$ is identified off the
extensive margin, and a single joint swarm spends most of its budget on the
former. Each sub-stage re-runs the swarm on one block, warm-started at the
incumbent and floored at its criterion value, so the criterion is monotone
across sub-stages by construction. The efficient fit of Step~3 restricts the
criterion to the moment blocks that carry the bootstrap covariance --- the
extensive-margin coefficients, the area-aggregated sourcing shares and the
count moments --- and re-optimises $\alpha$ within them, holding
$\{A_r,\Omega_L,\Omega_s\}$ at $\widehat{\mu}_1$; the remaining blocks are the
input--output and sales targets, which are treated as fixed and carry no
sampling variance.
 it where you want the
%  algorithms to appear (e.g. right after Appendix B.3.2).
%
%  Required packages in the master preamble:
%      \usepackage{amsmath,amssymb}
%      \usepackage{algorithm}
%      \usepackage{algpseudocode}
%      \usepackage{float}          % for the [H] float placement used below
%
%  Cross-references: equation numbers are written literally as (B12), (B22),
%  (B31), (B32) to match the current draft. Replace them by \eqref{...} once
%  those equations carry labels.
%  Recommended (renames Require/Ensure to Input/Output):
%      \renewcommand{\algorithmicrequire}{\textbf{Input:}}
%      \renewcommand{\algorithmicensure}{\textbf{Output:}}
%
%  `pseudocode.tex` in the same folder is a standalone driver that compiles
%  this file on its own, for checking.
% =====================================================================

\subsection{Pseudo-code}
\label{app:pseudocode}

This section states the two algorithms of Appendix~B.3 in
explicit form: Algorithm~\ref{alg:simulate} evaluates the model at a candidate
parameter vector $\mu$ and returns the simulated moment vector
$\widetilde{M}(\mu)$, and Algorithm~\ref{alg:smm} minimises the SMM criterion
over $\mu$ by particle swarm optimisation. Algorithms~\ref{alg:invertT},
\ref{alg:countN} and \ref{alg:pso} are the three subroutines they call: the
comparative-advantage inversion, the variety-count calibration, and the swarm
itself.

Throughout, $\mu=\bigl(\Omega_L,\{\Omega_s\},\alpha,\{A_r\}\bigr)$ is the vector
searched over; $a(r')$ denotes the attraction area of region $r'$;
$\mathcal{R}_a$ the regions of area $a$; $\mathcal{R}_D$ the set of downstream
buyers, $R_D=|\mathcal{R}_D|$; and $\mathcal{R}^+_s$ the regions inside active
attraction areas for sector $s$. The calibrated objects
\[
  \bigl\{\theta,\ \varepsilon,\ \lambda,\ \nu,\ \{\nu_s\},\
  \{w_{r's}\},\ \{\delta_r\},\ \{d_{r'r}\}\bigr\}
\]
are fixed once and for all and are suppressed from the argument lists.

% ---------------------------------------------------------------------
%  Algorithm 1 -- simulation
% ---------------------------------------------------------------------
\begin{algorithm}[H]
\caption{Simulation of the economy at a candidate $\mu$ (one criterion evaluation)}
\label{alg:simulate}
\begin{algorithmic}[1]
\Require candidate $\mu=(\Omega_L,\{\Omega_s\},\alpha,\{A_r\})$; empirical
  sourcing shares $\widehat{\gamma}_{as}$; empirical count moments
  $\widehat{G}_s(0)$ and bounds $[N^{lo}_s,N^{hi}_s]$; a fixed array of uniform
  draws $u_{r'\rho s}\in(0,1)$, $\rho=1,\dots,m$
\Ensure simulated moment vector $\widetilde{M}(\mu)$, and the calibrated
  $\{T_{as}\},\{\widehat{N}_s\}$
\vspace{2pt}
\State $\tau_{r'r}\gets d_{r'r}^{\,\alpha}$
  \Comment{trade costs; the only place $\alpha$ enters}
\vspace{2pt}
\State \textbf{Step 1 (comparative advantage).}
  $\{T_{as}\}\gets\Call{InvertT}{\alpha,\Omega_L,\{\Omega_s\},\{A_r\};\widehat{\gamma}}$
  \Comment{Algorithm~\ref{alg:invertT}}
\vspace{2pt}
\State \textbf{Step 2 (Ricardian competition).}
\For{$s=1,\dots,S$ \textbf{and} $\rho=1,\dots,m$}
    \For{$r'=1,\dots,R$}
        \State $z_{r'\rho s}\gets T_{a(r')s}^{1/\theta}\bigl(-\log u_{r'\rho s}\bigr)^{-1/\theta}$
          \Comment{Fr\'echet inverse c.d.f.}
    \EndFor
    \For{$r\in\mathcal{R}_D$}
        \State $r^\star(\rho,s,r)\gets\arg\min_{r'}\; w_{r's}\tau_{r'r}/z_{r'\rho s}$
          \Comment{winner of variety $\rho$ in market $(r,s)$}
        \State $p_{\rho s r}\gets w_{r^\star s}\tau_{r^\star r}/z_{r^\star\rho s}$
    \EndFor
\EndFor
\State $k_{r's}\gets\#\{\rho:\ r'=r^\star(\rho,s,r)\ \text{for some }r\in\mathcal{R}_D\}$,
  \quad $p_{r's}\gets k_{r's}/m$
  \Comment{win-anywhere counts and probabilities}
\vspace{2pt}
\State \textbf{Prices, sales and flows.} For every $r\in\mathcal{R}_D$:
\State \quad $P_{rs}\gets\bigl[\tfrac1m\sum_{\rho}p_{\rho s r}^{\,1-\nu_s}\bigr]^{1/(1-\nu_s)}$,\quad
  $P_r\gets\bigl[\sum_s\Omega_s P_{rs}^{1-\nu}\bigr]^{1/(1-\nu)}$
\State \quad $c_r\gets\bigl[\Omega_L w_r^{1-\lambda}+(1-\Omega_L)P_r^{1-\lambda}\bigr]^{1/(1-\lambda)}$,\quad
  $\tilde c_r\gets c_r/A_r$
\State \quad $y_r\gets \tilde c_r^{\,\varepsilon}\,\delta_r\big/\sum_{r''}\tilde c_{r''}^{\,\varepsilon}\delta_{r''}$,\quad
  $\pi_r\gets y_r$
  \Comment{iso-elastic demand, $E$ normalised to $1$}
\State \quad $X_{r'rs}\gets$ expenditure of $r$ on the varieties won by $r'$,
  from~(B12)
\vspace{2pt}
\State \textbf{Step 3 (variety count).}
  $\{\widehat{N}_s\}\gets\Call{CalibrateN}{\{k_{r's}\},m;\widehat{G}_s(0),[N^{lo}_s,N^{hi}_s]}$
  \Comment{Algorithm~\ref{alg:countN}}
\vspace{2pt}
\State \textbf{Step 4 (moments).} Assemble $\widetilde{M}(\mu)$ from the single
  solve above:
\State \quad (1) $\widetilde{\Omega}_L$, (2) $\widetilde{\Omega}_s$, (3) $\pi_r$
  \Comment{value blocks, from $\{X_{r'rs}\},\{c_r\},\{y_r\}$}
\State \quad (4) $\beta$: cloglog of ``does not supply'' on $\log d_{r'a(r')}$ and
  $\log z$, area${}\times{}$sector fixed effects, run on the simulated firm
  population $\{(r',\rho,s)\}$
\State \quad (5) $\widetilde{\gamma}_{as}\gets\sum_{r'\in\mathcal{R}_a}\sum_r X_{r'rs}\big/\sum_{r''}\sum_r X_{r''rs}$
\State \quad (6) $\widetilde{G}_s(\widehat{N}_s)$
  \Comment{count moment, evaluated at the calibrated $\widehat{N}_s$}
\State \Return $\widetilde{M}(\mu)$
\end{algorithmic}
\end{algorithm}

Three features of Algorithm~\ref{alg:simulate} are worth emphasising. First, the
arrows of~(B22) run one way: $T$ is calibrated without knowing
$N$, and $N$ does not revise $T$, so Steps~1--3 form a recursion rather
than a simultaneous system. Second, the simulation size $m$ is a Monte Carlo
device and is unrelated to $\widehat{N}_s$: the winner of a variety depends only
on relative production costs, so $p_{r's}$ can be estimated at any $m$ large
enough. Third, the uniform draws $u$ are held fixed across every evaluation of
the criterion within an optimisation stage (common random numbers), which is
what makes the criterion a deterministic function of $\mu$ that a swarm can
descend; they are scrambled Sobol sequences, constructed one low-dimensional
sequence per sector with regions as dimensions and varieties as points.
Independent pseudo-random draws are used instead when computing $\Sigma_{sim}$
and the numerical derivatives, where independence across replications is
required.

% ---------------------------------------------------------------------
%  Algorithm 2 -- comparative advantage inversion
% ---------------------------------------------------------------------
\begin{algorithm}[H]
\caption{\textsc{InvertT}: comparative advantage by matrix scaling with
general-equilibrium expenditure weights}
\label{alg:invertT}
\begin{algorithmic}[1]
\Require $\alpha,\Omega_L,\{\Omega_s\},\{A_r\}$; targets $\widehat{\gamma}_{as}$;
  damping $\varsigma\in(0,1]$; tolerance $\epsilon_T$; starting values $T^{(0)}_{as}$
\Ensure $\{T_{as}\}$ reproducing $\widehat{\gamma}_{as}$ exactly
\State $T\gets T^{(0)}$; normalise $T_{as}\gets T_{as}/T_{a_s^{\text{ref}}s}$ for each $s$
  \Comment{one reference area per sector fixes the scale}
\For{$t=1,2,\dots$ \textbf{until} convergence}
    \State $B_{ar}\gets\sum_{r'\in\mathcal{R}_a}(w_{r's}\tau_{r'r})^{-\theta}$,\quad
      $\Phi_{rs}\gets\sum_{a'}T_{a's}B_{a'r}$
    \State $\omega_{rs}\gets$ expenditure weights implied by $T$
      \Comment{price and sales block of Algorithm~\ref{alg:simulate}}
    \State $\gamma_{as}(T)\gets\sum_r\omega_{rs}\,T_{as}B_{ar}/\Phi_{rs}$
    \State $T^{+}_{as}\gets T_{as}\bigl(\widehat{\gamma}_{as}/\gamma_{as}(T)\bigr)^{\varsigma}$
      \Comment{multiplicative Sinkhorn update, in logs}
    \State $T^{+}_{as}\gets T^{+}_{as}/T^{+}_{a^{\text{ref}}_s s}$ for each $s$
    \State $\text{resid}\gets\max_{a,s}\bigl|\log T^{+}_{as}-\log T_{as}\bigr|$;\quad $T\gets T^{+}$
    \State \textbf{if} $\text{resid}<\epsilon_T$ \textbf{then break}
\EndFor
\State \Return $T$
\end{algorithmic}
\end{algorithm}

With expenditure weights held fixed, the multiplicative update is a Sinkhorn scaling of the strictly positive matrix
$B$, hence a contraction with a fixed point unique up to the sector-specific
scale removed by the reference-area normalisation. Recomputing $\omega_{rs}$
inside the loop adds the general-equilibrium response of prices and sales, so
the full map is a fixed-point iteration; the damping $\varsigma$ ($0.9$ in the
implementation) and the stopping rule on $\max_{a,s}|\Delta\log T_{as}|$ are what
the text refers to as requiring both sourcing shares and comparative advantages
to converge.

% ---------------------------------------------------------------------
%  Algorithm 3 -- variety count
% ---------------------------------------------------------------------
\begin{algorithm}[H]
\caption{\textsc{CalibrateN}: variety count by monotone integer search}
\label{alg:countN}
\begin{algorithmic}[1]
\Require win counts $k_{r's}$ from $m$ simulated varieties; targets
  $\widehat{G}_s(0)$; bounds $[N^{lo}_s,N^{hi}_s]$
\Ensure $\{\widehat{N}_s\}$
\For{$s=1,\dots,S$}
    \State define
      $\widetilde{G}_s(n)\gets\dfrac{1}{|\mathcal{R}^+_s|}\displaystyle\sum_{r'\in\mathcal{R}^+_s}
      \binom{m-k_{r's}}{n}\Big/\binom{m}{n}$
      \Comment{unbiased; decreasing in $n$}
    \If{$\widetilde{G}_s(N^{lo}_s)\le\widehat{G}_s(0)$}
        \State $\widehat{N}_s\gets N^{lo}_s$ \Comment{clamped at the lower bound}
    \ElsIf{$\widetilde{G}_s(N^{hi}_s)\ge\widehat{G}_s(0)$}
        \State $\widehat{N}_s\gets N^{hi}_s$ \Comment{clamped at the upper bound}
    \Else
        \State $(\underline{n},\overline{n})\gets(N^{lo}_s,N^{hi}_s)$
        \While{$\overline{n}-\underline{n}>1$}
            \State $n\gets\lfloor(\underline{n}+\overline{n})/2\rfloor$
            \State \textbf{if} $\widetilde{G}_s(n)>\widehat{G}_s(0)$ \textbf{then}
              $\underline{n}\gets n$ \textbf{else} $\overline{n}\gets n$
        \EndWhile
        \State $\widehat{N}_s\gets\arg\min_{n\in\{\underline{n},\overline{n}\}}
          \bigl|\widetilde{G}_s(n)-\widehat{G}_s(0)\bigr|$
    \EndIf
\EndFor
\State \Return $\{\widehat{N}_s\}$
\end{algorithmic}
\end{algorithm}

Because $\widetilde{G}_s$ is decreasing in $n$ and is evaluated in closed form
from the winning matrix of Step~2, the search costs $O(\log(N^{hi}_s-N^{lo}_s))$
evaluations and never re-simulates the Ricardian competition. The count is an
integer at every evaluation: it is never relaxed to the real line and never
rounded. A bound that binds is recorded and reported rather than silently
absorbed.

% ---------------------------------------------------------------------
%  Algorithm 4 -- PSO
% ---------------------------------------------------------------------
\begin{algorithm}[H]
\caption{\textsc{Pso}: particle swarm minimisation of $\mathcal{Q}$ over a box}
\label{alg:pso}
\begin{algorithmic}[1]
\Require objective $\mathcal{Q}$; box $[\ell,u]\subset\mathbb{R}^{d}$; swarm size
  $P$; iterations $I$; warm start $x_0$; inertia range
  $[\underline{w},\overline{w}]=[0.4,0.9]$; cognitive $c_1=1.5$; social $c_2=2.5$
\Ensure $x^\star=\arg\min \mathcal{Q}$ over the visited points
\State draw $x_1,\dots,x_{P-1}$ from a scrambled Sobol net on $[\ell,u]$;\;
  $x_P\gets\Pi_{[\ell,u]}(x_0)$
  \Comment{the warm start is always a particle}
\State $v_i\gets 0$; \; $b_i\gets x_i$, $f_i\gets\mathcal{Q}(x_i)$ evaluated in
  parallel; \; $g\gets\arg\min_i f_i$
\For{$t=1,\dots,I$}
    \State $w\gets\overline{w}-(\overline{w}-\underline{w})\,t/I$
      \Comment{inertia declines linearly}
    \For{$i=1,\dots,P$ \textbf{in parallel}}
        \State draw $r_1,r_2\sim U[0,1]^d$
        \State $v_i\gets w\,v_i+c_1 r_1\odot(b_i-x_i)+c_2 r_2\odot(g-x_i)$
        \State $v_i\gets\Pi_{[-\bar v,\bar v]}(v_i)$, \; $\bar v=0.1\,(u-\ell)$
          \Comment{velocity clamping}
        \State $x_i\gets x_i+v_i$; reflect $x_i$ back into $[\ell,u]$, damping the
          reflected velocity
        \State enforce the parameter constraints (monotone distance bins when
          $\alpha$ is binned)
        \State $f\gets\mathcal{Q}(x_i)$
          \Comment{one call to Algorithm~\ref{alg:simulate}}
        \State \textbf{if} $f<\mathcal{Q}(b_i)$ \textbf{then} $b_i\gets x_i$;\;
          \textbf{if} $f<\mathcal{Q}(g)$ \textbf{then} $g\gets x_i$
    \EndFor
\EndFor
\State \Return $g$
\end{algorithmic}
\end{algorithm}

% ---------------------------------------------------------------------
%  Algorithm 5 -- estimation
% ---------------------------------------------------------------------
\begin{algorithm}[H]
\caption{SMM estimation of $\mu$ by particle swarm optimisation}
\label{alg:smm}
\begin{algorithmic}[1]
\Require empirical moments $\widehat{M}$; empirical sourcing shares
  $\widehat{\gamma}_{as}$ and counts $\widehat{G}_s(0)$; bootstrap covariance
  $\Sigma_{data}$; number of refinement loops $L$; swarm size $P$; iterations $I$
\Ensure $\widehat{\mu}_2$, standard errors, and the calibrated
  $\widehat{T},\widehat{N}$ with their standard errors
\vspace{2pt}
\State \textbf{Starting values.} $\Omega_L,\{\Omega_s\}$ at their input--output
  counterparts; $A_r\propto\pi_r^{1/\varepsilon}$; $\alpha$ at the
  extensive-margin coefficient divided by $\theta$. Call this $\mu^{(0)}$, and
  set the box $[\ell,u]$ multiplicatively around it.
\vspace{3pt}
\State \textbf{Step 1: first fit under equal weights.} Let
  $\mathcal{Q}_I(\mu)=\bigl(\widetilde{M}(\mu)-\widehat{M}\bigr)'\bigl(\widetilde{M}(\mu)-\widehat{M}\bigr)$,
  each evaluation running Algorithm~\ref{alg:simulate}.
\State \quad $\mu\gets\Call{Pso}{\mathcal{Q}_I,[\ell,u],P,I,\mu^{(0)}}$
  \Comment{joint fit over all of $\mu$}
\For{$\text{loop}=1,\dots,L$}
    \Comment{alternating block refinement, search radius shrinking with the loop}
    \For{$\mathcal{B}\in\bigl(\{A_r\},\ \{\alpha\},\ \{\Omega_L,\Omega_s\}\bigr)$}
        \State $\mu\gets\Call{Pso}{\mathcal{Q}_I,\mathcal{B}\text{-box of radius }
          \varrho_{\text{loop}}\text{ around }\mu,P,I,\mu}$
        \Comment{coordinates outside $\mathcal{B}$ fixed}
    \EndFor
    \State \textbf{break if} $\max_k|\Delta\mu_k|/|\mu_k|<\epsilon_\mu$
\EndFor
\State $\widehat{\mu}_1\gets\mu$
\vspace{3pt}
\State \textbf{Step 2: the weighting matrix.} At $\widehat{\mu}_1$, re-simulate
  $K$ times with independent draws, $\Sigma_{sim}\gets\widehat{\mathrm{Var}}
  \bigl(\widetilde{M}^{(1)},\dots,\widetilde{M}^{(K)}\bigr)$, and set
  $W\gets(\Sigma_{data}+\Sigma_{sim})^{-1}$.
\vspace{3pt}
\State \textbf{Step 3: the efficient fit.} Let
  $\mathcal{Q}_W(\mu)=\bigl(\widetilde{M}(\mu)-\widehat{M}\bigr)'W
  \bigl(\widetilde{M}(\mu)-\widehat{M}\bigr)$ and repeat the joint fit and the block
  refinement of Step~1, starting from
  $\widehat{\mu}_1$, yielding $\widehat{\mu}_2$.
\vspace{3pt}
\State \textbf{Step 4: inference.} At $\widehat{\mu}_2$:
\State \quad $G\gets\partial\widetilde{M}/\partial\mu$ by central finite
  differences with common random numbers
\State \quad report $(G'WG)^{-1}$ and
  $(G'WG)^{-1}G'W\Omega WG(G'WG)^{-1}$, $\Omega=\Sigma_{data}+\Sigma_{sim}$
\State \quad $\widehat{T}$: delta method on the calibration map
  $T^\star(\alpha,\widetilde{\gamma})$, equation~(B31)
\State \quad $\widehat{N}$: delta method on
  $\widetilde{G}_s(\widehat{N}_s)=\widehat{G}_s(0)$, equation~(B32)
\State \Return $\widehat{\mu}_2$ and the associated standard errors
\end{algorithmic}
\end{algorithm}

Two implementation notes. The block refinement inside Step~1 exists because the
blocks of $\mu$ have very different curvature: $\{A_r\}$ is pinned almost
one-for-one by the regional sales shares, while $\alpha$ is identified off the
extensive margin, and a single joint swarm spends most of its budget on the
former. Each sub-stage re-runs the swarm on one block, warm-started at the
incumbent and floored at its criterion value, so the criterion is monotone
across sub-stages by construction. The efficient fit of Step~3 restricts the
criterion to the moment blocks that carry the bootstrap covariance --- the
extensive-margin coefficients, the area-aggregated sourcing shares and the
count moments --- and re-optimises $\alpha$ within them, holding
$\{A_r,\Omega_L,\Omega_s\}$ at $\widehat{\mu}_1$; the remaining blocks are the
input--output and sales targets, which are treated as fixed and carry no
sampling variance.
